\chapter{Ejecucion, seguridad y cumplimiento normativo}
\thispagestyle{memoria}\pagestyle{memoria}

\section{Secuencia de ejecucion}
La instalacion se ejecutara por frentes coordinados con arquitectura y con la
especialidad de hidrocarburos. Primero se replantean rutas, tableros, cruces y
limites de areas clasificadas; luego se instalan canalizaciones y cajas; despues
se tienden, identifican y prueban conductores; finalmente se montan tableros,
equipos, alumbrado, puesta a tierra, respaldo y controles. No se energizara una
zona de combustible mientras existan trabajos con llama, empalmes provisionales
o equipos sin certificacion aplicable.

\section{Criterios de valorizacion}
Las partidas de canalizacion y cableado se valorizan por metro instalado, con
10\% de holgura para despieces, ascensos, reservas y desperdicio controlado. Las
luminarias, tableros, protecciones, UPS, ATS, electrodos y puntos se valorizan por
unidad o conjunto probado. Los equipos propios del proceso de combustible
(tanques, surtidores y STP) son suministro de otra especialidad: este presupuesto
solo incluye su alimentacion, proteccion, control y conexion equipotencial.

No se aceptan metrados duplicados por ``punto completo'' cuando la canalizacion y
los conductores ya aparecen por metro. Toda valorizacion exige protocolo de
prueba, rotulo y correspondencia con el circuito del plano IE-05.

\section{Normas base}
El diseño se apoya en el Codigo Nacional de Electricidad--Utilizacion aprobado por RM 037-2006-MEM/DM y sus reglas aplicables, la Norma EM.010 del RNE aprobada por RM 083-2019-VIVIENDA, el DS 054-93-EM y sus modificatorias, y la Guia Tecnica 001 de OSINERGMIN como apoyo facultativo. La lista exacta, fecha de consulta y matriz de cumplimiento se conservan en la documentacion del repositorio.

No se afirma una equivalencia automatica entre las Zonas 0/1/2 del CNE-U y las divisiones empleadas por reglas sectoriales. La lamina IE-06 es una propuesta academica que debe ser revisada por especialista competente con equipos, alturas, ventilacion y proceso definitivos.

\section{Areas peligrosas}
Se considera Zona 0 en interiores de tanques y tuberias con presencia continua o prolongada de vapor; Zona 1 en las envolventes inmediatas de llenado/venteo definidas por CNE-U 120; y Zona 2 alrededor de surtidores y envolventes secundarias. En planta se representa el radio horizontal de 6 m alrededor de cada surtidor y las envolventes adoptadas de bocas de llenado y venteos.

La clasificacion real cambia si varian producto, caudal, altura, ventilacion, drenaje, equipos o geometria. Por ello los planos llevan advertencia de anteproyecto y no autorizan construccion.

\section{Pendientes obligatorios antes de obra}
\begin{itemize}
  \item Factibilidad, punto de entrega, potencia e Icc de Electro Puno.
  \item Placas y fichas de surtidores, STP, compresor, bombas, refrigeracion, UPS y grupo.
  \item Levantamiento acotado, altura de venteos/llenado y verificacion de las unidades del DXF.
  \item Resistividad de suelo, medicion de PAT y evaluacion de riesgo de rayo.
  \item Calculo de cortocircuito, selectividad y coordinacion de protecciones.
  \item Simulacion fotometrica, estudio de ventilacion y revision profesional de areas clasificadas.
  \item Firma, sello y CIP de profesional habilitado si el trabajo se convierte en expediente constructivo.
\end{itemize}
